\scanpage{108}

\uline{Genre d'une fonction entière}

Il arrive que l'on ait $\sum\frac{1}{|a_n|^{k+1}}<+\infty$ pour quelque $k$, (soit $k$ le plus petit entier qui y satisfasse) et de plus que dans la repré-\newline sentation (2), $g(z)$ soit un polynôme -- d'ailleurs déterminé à une constante près.

Soit $r$ son degré, le plus \deleted{petit} \inserted{grand} des nombres $r,k$ s'appelle le genre de $f(z)$.
\inserted{C'est un nb \uline{entier} $\geq 0$.}

Donc si $f(z)$ est de genre $\leq r$, alors on a la formule
\[
f(z)=e^{\alpha_0+\alpha_1 z\mathord{\cdot}\mathord{-}+\alpha_n z^n}\prod_n\left(1-\frac{z}{a_n}\right)e^{\frac{z}{a_n}+\mathord{-}\mathord{\cdot}\frac{z^{\unsure{n|r}}}{r a_n}}\tag*{(3)}
\]

Dire que $f(z)$ est de genre zéro, signifie que $f(z)$ est de la forme
\[
\begin{aligned}\vphantom{f(z)}\\[-2pt]\text{avec}\quad\end{aligned}
\left\{\begin{aligned}
f(z)&=\mathrm c^{\mathrm{te}}\times\prod_n\left(1-\frac{z}{a_n}\right)\unsure{.|,}\\[-2pt]
\sum_n\frac{1}{|a_n|}&<+\infty
\end{aligned}\right.\tag*{(4)}
\]

\uline{Ordre d'une fonction entière} \quad Posons $M(r)=\mathop{\rm Sup}_{z\unsure{=|\leq}r}|f(z)|$

Puis $\rho=\overline{\lim}\frac{\log\log M(r)}{\substack{\deleted{\unsure{z}}\\\inserted{\log r}}}$ \quad $(0\leq\rho\leq+\infty)$, $\rho$ est dit l'ordre de $f(z)$, et $f(z)$ est dite d'ordre
\par\noindent\hspace{2cm}\unsure{$=$}\par\noindent
fini si $\rho<+\infty$. Cela signifie qu'il existe $k>0$ tel que $|f(z)|\leq e^{\unsure{(z\vert^k)|(\lvert z\rvert^k)}}$. Alors $\rho$ est la \deleted{\unsure{plus petit nombre|plus petite valeur}} \inserted{borne inférieure des} $\sigma>0$ tels que

\margin{109}
